\documentclass{article}
\usepackage[T1]{fontenc}
\usepackage{lmodern}
\usepackage{graphicx}
\usepackage{amsmath,amssymb}
\usepackage[a4paper,margin=2cm]{geometry}
\usepackage[absolute,overlay]{textpos}
\setlength{\TPHorizModule}{1mm}
\setlength{\TPVertModule}{1mm}
\usepackage[indonesian]{babel}
\usepackage{hyperref}
\setlength{\emergencystretch}{2em}
% unit: Halaman 14

\begin{document}

% === Halaman 14 (llm) ===

\section*{RC5}

\begin{itemize}
    \item $RC5$ dibuat oleh Ronald Rivest dari Laboratorium $RSA$ pada tahun 1994.
    \item Salah satu finalis AES, yaitu RC6, dikembangkan dari RC5
\end{itemize}

\begin{itemize}
    \item Tidak seperti algoritma \textit{cipher} blok lainnya, $RC5$ mempunyai:
    \begin{itemize}
        \item ukuran blok yang variabel (16, 32, 64)
        \item panjang kunci yang variabel (0 sampai 2040 bit)
        \item dan jumlah putaran yang variabel (0 sampai 255).
    \end{itemize}
\end{itemize}

\vspace{5mm}

\begin{tabular}{|l|c|l|}
\hline
\textbf{Parameter} & \textbf{Simbol} & \textbf{Nilai yang dibolehkan} \\ \hline
Ukuran blok (dalam bit) & $w$ & 16, 32, 64 \\ \hline
Jumlah putaran & $r$ & 0, 1, ..., 255 \\ \hline
Panjang kunci eksternal $K$ (dalam $byte$, 1 $byte = 8$ bit) & $b$ & 0, 1, ..., 255 \\ \hline
\end{tabular}

\end{document}
